% ============================================================
%  CHAPTER 3: METHODOLOGY / SYSTEM DESIGN
% ============================================================
%
%  VAI TRO CUA CHAPTER 3:
%  Day la TRAI TIM cua luan van ky thuat. Ban mo ta THIET KE
%  cua he thong / phuong phap: tai sao chon kien truc nay,
%  tai sao thiet ke tung component nhu vay.
%
%  PHAN BIET "METHODOLOGY" VA "IMPLEMENTATION":
%    METHODOLOGY (Chapter 3) = THIET KE + LY THUYET
%      "What the system does and why it was designed that way."
%      Tap trung vao kien truc, quyet dinh thiet ke, cac cong
%      thuc va ly thuyet dung lam nen tang.
%
%    IMPLEMENTATION (Chapter 4) = HIEN THUC HOA
%      "How the system was built in practice."
%      Tap trung vao cong cu, thu vien, code, cau hinh cu the.
%
%  VID DU PHAN BIET:
%    Chapter 3: "The system uses cosine similarity to rank
%                retrieved documents because it is invariant
%                to vector magnitude, unlike Euclidean
%                distance, making it more robust for comparing
%                sentence embeddings of different lengths."
%    Chapter 4: "Cosine similarity is computed using the
%                Qdrant client library (v1.7.0) with
%                HNSW indexing, configured at ef=128."
%
%  LOI SAI PHO BIEN:
%  1. Viet Chapter 3 nhu huong dan cai dat (setup guide).
%     Chapter 3 phai giai thich QUYET DINH, khong chi HANH DONG.
%  2. Mo ta HET MOI THU ma khong giai thich tai sao.
%     Moi component phai co ly do chon no.
%  3. Quen khong co system overview / architecture figure.
%     Phai co it nhat 1 hinh mo ta toan bo pipeline.
% ============================================================

\chapter{Methodology}
\label{chap:methodology}

% ============================================================
%  OPENING PARAGRAPH (bat buoc)
%  -----------------------------------------------------------
%  Gioi thieu ngan gon cau truc cua chapter nay.
%  Neu ro cac section chinh theo thu tu pipeline / thiet ke.
%
%  Vi du:
%  "This chapter presents the design of the proposed [system].
%   Section 3.1 provides a high-level overview of the
%   system architecture. Sections 3.2 through 3.5 then
%   describe each major component in the order they appear
%   in the processing pipeline: [component 1], [component 2],
%   [component 3], and [component 4].
%   Section 3.6 explains the parameter selection rationale."
%
%  LUU Y: Thu tu cac section nen theo thu tu xu ly cua he thong.
%  Dieu nay giup nguoi doc hieu luong du lieu di qua dau.
% ============================================================

% *** VIET OPENING PARAGRAPH TAI DAY ***
[Write the opening paragraph that maps out this chapter.
Follow the pipeline order -- input -> processing -> output.]


% ============================================================
%  SECTION 3.1: SYSTEM OVERVIEW
%  -----------------------------------------------------------
%  Day la section quan trong nhat trong Chapter 3.
%  Nguoi doc can thay "buc tranh lon" truoc khi di vao chi tiet.
%
%  NOI DUNG BAT BUOC:
%  (1) Mo ta tong the he thong lam gi
%  (2) Liet ke cac component chinh
%  (3) Mo ta luong xu ly (processing flow) tu dau den cuoi
%  (4) FIGURE: So do kien truc he thong (architecture diagram)
%
%  HUONG DAN VE FIGURE KIEN TRUC:
%  - Phai co it nhat 1 figure o Section 3.1
%  - Figure phai hien thi: input -> cac component -> output
%  - Ten cac component trong figure phai khop voi ten trong text
%  - Figure caption phai mo ta du (ai khong doc text van hieu)
%
%  VI DU VIET SYSTEM OVERVIEW TOT:
%  "The system operates as a pipeline with N stages.
%   When a query arrives, it is first [step 1: preprocessed].
%   [Step 2: Intent classification] then determines which
%   retrieval path the query should follow. Depending on the
%   classified intent, the query is routed to either
%   [Path A] or [Path B]. The retrieved content is then
%   passed to [generation component], which produces the
%   final response. Figure~\ref{fig:arch} illustrates this
%   flow."
%
%  MIEN DICH TU VAN DE THUC TE SANG KIEN TRUC:
%    Van de: "Students ask prerequisite questions"
%    -> Component: "Intent classifier labels query as
%                   'prerequisite'"
%    -> Action: "System queries SQL database"
%    -> Output: "Templated answer returned"
%  Lam ro khi nao dung rule-based, khi nao dung ML.
% ============================================================

\section{System Overview}
\label{sec:overview}

% *** VIET SYSTEM OVERVIEW TAI DAY ***
[Describe the overall system in 2-4 paragraphs. Cover:
(1) The problem the system solves at a high level
(2) The main components (list them by name)
(3) The processing flow from input to output
(4) Reference the architecture figure below]

\begin{figure}[H]
  \centering
  % CHEN HINH KIEN TRUC CUA BAN VAO DAY
  % Doi file hinh thanh dang PNG hoac PDF, luu vao thu muc images/
  \IfFileExists{images/architecture}{%
    \includegraphics[width=0.95\textwidth]{images/architecture}%
  }{%
    \fbox{\parbox[c][5cm][c]{0.9\textwidth}{\centering Architecture figure not added yet}}%
  }
  \caption{[System architecture overview: describe the main
    components and how data flows through them. A good caption
    should be self-contained -- a reader who sees only this
    figure and its caption should understand the overall flow.]}
  \label{fig:arch}
\end{figure}


% ============================================================
%  SECTION 3.2: [COMPONENT 1 -- e.g., Data / Knowledge Base]
%  -----------------------------------------------------------
%  Bat dau voi component dau tien trong pipeline.
%  Thuong la du lieu / knowledge base / bieu dien du lieu.
%
%  CAC DIEU CAN CO TRONG MOI COMPONENT SECTION:
%  (1) Mo ta component lam gi (1 doan)
%  (2) Ly do chon thiet ke nay (1 doan)
%      "We chose X over Y because..."
%      "This design decision was motivated by..."
%  (3) Chi tiet ky thuat (subsections neu phuc tap)
%  (4) Figure hoac table neu giup hieu ro hon
%
%  CHO PHAN DATA/KNOWLEDGE BASE, NEN BAO GOM:
%  - Nguon du lieu (data sources)
%  - Kich thuoc va dac diem du lieu
%  - Cach xu ly / clean / format du lieu
%  - Bieu dien du lieu (representation): tai sao chon format nay
% ============================================================

\section{[Component 1: e.g., Knowledge Base Construction]}
\label{sec:knowledge_base}

% *** VIET NOI DUNG COMPONENT 1 TAI DAY ***

[Paragraph 1: What does this component do and what problem
does it solve?]

[Paragraph 2: What are the data sources? Describe size,
format, and characteristics.]

[Paragraph 3: How did you process or structure the data?
Why did you choose this representation?]

% Chi them subsections neu component nay phuc tap
% \subsection{Data Sources}
% \subsection{Data Preprocessing}
% \subsection{Document Representation}


% ============================================================
%  SECTION 3.3: [COMPONENT 2 -- e.g., Core Technical Method]
%  -----------------------------------------------------------
%  Day thuong la phan ky thuat chinh cua he thong.
%  Can co CONG THUC TOAN HOC neu phuong phap co nen ly thuyet.
%
%  HUONG DAN VE CONG THUC:
%  - Moi cong thuc can duoc danh so: \begin{equation}...\end{equation}
%  - Giai thich tat ca ky hieu ngay sau cong thuc:
%    "where $\mathbf{v}_q$ is the query vector, $k$ is..."
%  - Giai thich TAI SAO dung cong thuc nay, khong chi la no la gi
%
%  Vi du viet cong thuc tot:
%  "Similarity between query and document is computed as:
%   \begin{equation}
%     \text{sim}(\mathbf{v}_q, \mathbf{v}_d) = 
%     \frac{\mathbf{v}_q \cdot \mathbf{v}_d}
%          {\|\mathbf{v}_q\| \cdot \|\mathbf{v}_d\|}
%   \end{equation}
%   where $\mathbf{v}_q$ is the query embedding vector and
%   $\mathbf{v}_d$ is the document embedding vector.
%   Cosine similarity was chosen over Euclidean distance
%   because it is invariant to vector magnitude, which
%   matters when comparing sentence embeddings of varying
%   lengths."
% ============================================================

\section{[Component 2: e.g., Retrieval / Core Method]}
\label{sec:retrieval}

% *** VIET NOI DUNG COMPONENT 2 TAI DAY ***

[Paragraph 1: Describe what this component does in the
overall pipeline.]

[Paragraph 2: Describe the core technical mechanism.
Include the key formula if applicable.]

% Example formula -- replace with your own:
\begin{equation}
  \text{[formula name]} = \frac{[numerator]}{[denominator]}
  \label{eq:main}
\end{equation}

where $[\text{symbol 1}]$ is [description], and
$[\text{symbol 2}]$ is [description].

[Paragraph 3: Justify your design choices. Why this
approach and not an alternative?]


% ============================================================
%  SECTION 3.4: [COMPONENT 3]
%  -----------------------------------------------------------
%  Tiep tuc moi thanh phan trong pipeline.
%  Lap lai cau truc: mo ta -> ly do chon -> chi tiet ky thuat
% ============================================================

\section{[Component 3: e.g., Classification / Routing]}
\label{sec:classification}

% *** VIET NOI DUNG COMPONENT 3 TAI DAY ***

[Describe this component following the same pattern:
(1) What it does, (2) Why this design, (3) Technical details]

% Include a table if this component involves multiple cases:
% Example: Intent categories and their handling paths
\begin{table}[H]
  \centering
  \caption{[Table title -- describe what the table shows]}
  \label{tab:component3}
  \begin{tabular}{lp{5cm}p{4cm}}
    \toprule
    \textbf{{[Column 1]}} & \textbf{{[Column 2]}} & \textbf{{[Column 3]}}\\
    \midrule
    {[row 1, col 1]} & {[row 1, col 2]} & {[row 1, col 3]}\\
    {[row 2, col 1]} & {[row 2, col 2]} & {[row 2, col 3]}\\
    \bottomrule
  \end{tabular}
\end{table}


% ============================================================
%  SECTION 3.5: [COMPONENT 4 -- e.g., Generation / Output]
%  -----------------------------------------------------------
%  Component cuoi cung trong pipeline: tao ra output.
%  Can mo ta ro rang: khi nao dung rule-based template,
%  khi nao dung language model / ML model.
% ============================================================

\section{[Component 4: e.g., Response Generation]}
\label{sec:generation}

% *** VIET NOI DUNG COMPONENT 4 TAI DAY ***

[Describe the output generation component.
Be explicit about: when is the output deterministic/templated
vs. when does it use a learned model? Why?]


% ============================================================
%  SECTION 3.6: PARAMETER SELECTION (neu co)
%  -----------------------------------------------------------
%  Neu he thong co cac tham so quan trong (k, threshold,
%  model size, etc.), giai thich tai sao ban chon gia tri do.
%
%  Cach viet tot:
%  "The retrieval depth k=3 was determined empirically:
%   k=1 was insufficient when the top result was marginally
%   incorrect; k=5 and above diluted the context window.
%   k=3 provided the best balance in preliminary runs."
%
%  Neu ban toi uu tham so bang grid search hoac ablation,
%  tham chieu sang Chapter 4/5 noi ban bao cao ket qua.
% ============================================================

\section{Parameter Selection}
\label{sec:params}

% *** VIET PARAMETER SELECTION TAI DAY (hoac xoa section neu khong can) ***

[Explain the key parameters in your system and how you
selected their values. Was it empirical? Theoretical?
Based on prior work? Be honest about what was tested
vs. what was set by convention.]

